\documentclass[activite]{thedocument}
\startdatakeys
\source{}
\cycle{}
\competencesDuSocle{}
\connaissancesRequises{}
\competencesTravaillees{}
\enddatakeys
\begin{document}

\begin{enumerate}
    \item Associer à chaque égalité un (ou deux) schéma(s) correspondant(s).\vskip10pt
    \begin{minipage}{.5\linewidth}
        \begin{itemize}
            \item $20=4\times5$
            \item $5=1\times5$
            \item $12=3\times4$
            \item $9=3\times3$
        \end{itemize}
    \end{minipage}%
    \begin{minipage}{.5\linewidth}
        \pdf[scale=0.3]{activite-multiples_diviseurs_schema}
    \end{minipage}
    \item Dessiner le schéma correspondant à l'égalité suivantes~:
    \par$16=2\times8$\par\vskip100pt
    \item Peut-on faire d'autre schémas en utilisant une égalité du type
    $16=a\times b$ où $a$ et $b$ sont des nombres entiers~?\vskip50pt
    \item Dans ce cas, comment appelle-t-on les nombres $a$ et $b$ ?
    \dotedline\dotedline\dotedline
\end{enumerate}

\end{document}